\scanpage{23}
\mspage{4}

Posons
\[
\begin{aligned}M_0=\mathrm{End}_X(e)=\{\,f\in M\mid ef&=\deleted{\unsure{f}}e\ \text{i.e.}\ a\psi f=a\psi\\ &\text{i.e.}\ \unsure{a\psi|\psi}f\overset{R}\sim\psi\,\}\end{aligned}
\]

Ainsi, on a pour $f\in M$
\[
f\in M_0\quad\Longleftrightarrow\quad\forall\,p,q\in\mathbb N\qquad\varphi^p\psi f=\varphi^q\psi.
\]

Je vais écrire une condition sur $\varphi,\psi$ qui impliquera que pour $f\in M_0$ fixé, l'entier $\delta(f)\overset{\mathrm{def}}{=}q-p$ ne dépend pas du choix de $p,q$.

Supposons donc qu'on ait
\[
\left\{\begin{aligned}\varphi^p\psi f&=\varphi^q\psi\\ \varphi^{p'}\psi f&=\varphi^{q'}\psi\end{aligned}\right.
\]
on en conclut
\[
\varphi^{p+p'}\psi f=\left\{\begin{aligned}&\varphi^{p'+q}\psi\\ &\varphi^{p+q'}\psi\end{aligned}\right.
\]
donc
% letterhead
\par\smallskip
{\footnotesize\noindent\begin{minipage}[t]{.46\linewidth}\raggedright
MINISTERE DE L'EDUCATION NATIONALE\\ ACADEMIE D'AIX-MARSEILLE\\ \rule{2.6em}{.3pt}\\ Inspection Académique\\ de Vaucluse\\ \rule{2.6em}{.3pt}\\ Tél. : 90 82 99 76\\ Référence : Secrétariat n\textsuperscript{o}
\end{minipage}\hfill%
\begin{minipage}[t]{.51\linewidth}\raggedright
AVIGNON, le\\ L'Inspecteur d'académie,\\ directeur des services départementaux\\ de l'Education Nationale de Vaucluse,\\ à M
\end{minipage}\par}

% archive sheet number: 22
